\newgeometry{margin=1in, twoside=false}
\setlength{\LTleft}{-2cm}
\begin{longtable}{>{\raggedright\arraybackslash}p{3cm}|>{\raggedright\arraybackslash}p{16cm}}
\caption{EvalOps Framework Capabilities Checklist: Evaluative definitions for assessing evaluation harness capabilities across 4 stages and 28 capabilities. These definitions guide systematic assessment of how well harnesses operationalize construct-valid evaluation.} \label{tab:evalops-checklist} \\
\hline
\textbf{Capability} & \textbf{Definition} \\
\hline

\multicolumn{2}{l}{\textbf{STAGE 1: CONFIGURE -- Evaluation Setup \& Specification (8 capabilities)}} \\
\hline

S1F1: Benchmark Loading \& Validation & This capability assesses whether the harness provides capabilities for loading construct-valid benchmarks as foundational input. A construct-valid benchmark consists of five components: a construct definition describing the phenomenon being measured, a pre-sampled dataset of evaluation instances, prompt specifications detailing task presentation, a measurement protocol explaining scoring approaches, and validity evidence documenting sampling rationale and limitations. Look for support for multiple benchmark formats such as JSON, YAML, CSV, or dataset repositories, such as HuggingFace. The harness should validate that required benchmark components are present and provide clear error messages when components are missing. Integrity verification through checksums or version tracking helps detect data corruption or drift. Support for loading multiple benchmarks simultaneously enables multi-task assessment. The harness treats benchmarks as complete packages where construct validity has already been established by benchmark designers, focusing on operationalizing rather than establishing that validity. \\
\hline

S1F2: System Specification \& Configuration & This capability examines what configuration options the harness provides for specifying the target system being evaluated. Configuration should cover four key dimensions: the model function with its parameters, version identifiers such as model IDs or checkpoint locations for exact reproducibility, inference settings controlling output generation including temperature and sampling parameters, and resource specifications defining compute allocation and operational constraints, such as memory limits and timeouts. Look for support across multiple model providers such as OpenAI, Anthropic, HuggingFace, local models, and custom endpoints. The harness should enable precise version pinning and provide validation of parameter values against model constraints. Support for configuring multiple systems simultaneously enables head-to-head comparisons, ablation studies, and baseline evaluations under identical conditions. Configuration should be serializable and reusable to support reproducible re-runs. \\
\hline

S1F3: Protocol Selection \& Configuration & This capability evaluates what measurement modalities the harness supports for scoring model outputs. The harness should provide interfaces for diverse measurement approaches. Human judgment support includes integration with crowdsourcing platforms, rater qualification mechanisms, attention checks, inter-rater reliability computation, and demographic documentation. LLM-as-judge support includes judge model selection, rubric design, chain-of-thought reasoning prompts, confidence calibration, and bias mitigation strategies. Algorithmic metrics support includes diverse scoring functions, such as BLEU, ROUGE, and semantic similarity with configurable parameters. Ensemble approaches combine multiple modalities through voting schemes and confidence weighting. Look for a standardized measurement interface where all modalities return scores, uncertainty estimates, rationales, and metadata. The harness should support custom metric definition while preserving aggregation of subjective judgments without forcing single-point estimates. \\
\hline

S1F4: Baseline Specification \& Configuration & This capability assesses what baseline system support the harness provides for establishing comparative context. Look for capabilities to configure diverse baseline types: random baselines using uniform random or majority-class predictions, classical methods, such as logistic regression and support vector machines with hyperparameter tuning, state-of-the-art models from published literature, and human performance benchmarks. The harness should provide mechanisms to ensure fair comparison conditions including identical data splits across all systems, consistent measurement protocols with the same raters or judges, equal hyperparameter tuning budgets, and comparable computational resources. Baseline results should be presented alongside target system performance to enable convergent validity evidence showing the target outperforms simpler alternatives, and discriminant validity assessment positioning performance relative to human capability. \\
\hline

S1F5: Resource Budget Planning & This capability assesses whether the harness provides cost estimation and budget management capabilities. Look for support to specify compute budgets in GPU/CPU hours, cost limits for API tokens and cloud compute, and time constraints for wall-clock deadlines and throughput requirements. The harness should estimate total costs before execution given current configuration, ideally using token-based modeling that accounts for input tokens, output tokens, sample count, and per-token pricing. Cost projections should inform tradeoff decisions about sample sizes, measurement protocol selection, baseline scope, and cross-validation fold counts. The harness should provide budget enforcement mechanisms that prevent runs from exceeding limits or alert when approaching thresholds. Cost reporting in results should break down expenses by component and compare actual versus estimated costs. \\
\hline

S1F6: Cross-Validation Configuration & This capability evaluates what cross-validation options the harness supports for structuring evaluation before model exposure to data. Look for support for multiple CV approaches: k-fold for general use, stratified k-fold for preserving class distributions, leave-one-out for small datasets, time-series forward chaining respecting temporal ordering, and nested CV for simultaneous tuning and evaluation. The harness should provide deterministic split generation through fixed random seeds or explicit partition indices, enabling identical split reproduction across runs. Stratification controls should allow balanced distributions across critical attributes. The harness should enforce test set isolation to prevent data leakage. Split specifications should be serializable for reproducibility, and CV structure should integrate properly with statistical aggregation to reflect per-fold variability in uncertainty estimates. \\
\hline

S1F7: Analysis Plan Specification & This capability examines whether the harness provides mechanisms for pre-specifying statistical analysis plans before evaluation execution. Look for support to define sample size requirements with power analysis capabilities, designation of primary versus secondary metrics, significance threshold specification, multiple comparison correction options, such as Bonferroni or Benjamini-Hochberg FDR, uncertainty quantification method selection such as bootstrap or parametric confidence intervals, and stratification variable identification. The harness should store and version these analysis plans, making it possible to verify that reported analyses follow pre-specified protocols rather than post-hoc explorations. Plan adherence helps prevent selective reporting and p-hacking by establishing analysis commitments before seeing results. The interface for specifying plans should be clear and the stored plans should be accessible for verification. \\
\hline

S1F8: Provenance Tracking \& Configuration & This capability examines what provenance tracking capabilities the harness provides for comprehensive reproducibility. Look for automatic capture of runtime metadata including execution timestamps and durations, configuration snapshots with all parameters serialized and defaults explicit, model identification through IDs and version hashes, environment specifications documenting OS and library versions with complete dependency trees, data provenance with cryptographic hashes and preprocessing pipelines, execution logs with errors and performance metrics, and git integration tracking configuration changes. The harness should allow configuration of metadata tracking scope and logging verbosity. Artifact storage should include compression options and integrity verification through checksums. Complete provenance enables bit-by-bit reproducibility where independent researchers can regenerate exact results, supports external validation and regulatory compliance, and facilitates debugging and longitudinal tracking. \\
\hline

\multicolumn{2}{l}{\textbf{STAGE 2: EXECUTE -- Operational Evaluation Execution (7 capabilities)}} \\
\hline

S2F1: Data Loading \& Preprocessing & This capability evaluates the harness's data handling capabilities across diverse sources and modalities. Look for support for multiple formats including structured formats, such as JSON and CSV, database queries, dataset repositories, such as HuggingFace, REST APIs, and cloud storage. Multi-modal support should handle text in various formats, images, audio, video, structured data, and combinations. The harness should apply preprocessing transformations specified by benchmarks including tokenization, normalization, augmentation, and format conversion. For large datasets exceeding memory, streaming support through batching should be available. Look for automatic schema standardization ensuring consistency across source formats, data validation checking accessibility and version consistency, and deterministic dataset splitting for reproducibility. Preprocessing configuration should be documented and allow customization or disabling of transformations. \\
\hline

S2F2: Schema Validation \& Consistency & This capability assesses what schema enforcement capabilities the harness provides before expensive inference begins. Look for type checking validating field data types, structure validation confirming required fields and correct nesting, range and bound checking verifying numerical values and categorical constraints, null and missing value handling with documented acceptable patterns, and cardinality constraint validation. The harness should verify training-serving consistency by ensuring the same preprocessing pipeline, capability engineering, and data representation used during model training. Schema versioning should track evolution across benchmark versions and detect backward compatibility issues. Distribution shift detection should flag schema drift that could invalidate comparisons or violate benchmark assumptions. Early validation catches data quality issues before consuming resources on invalid runs. \\
\hline

S2F3: Contamination Detection & This capability evaluates what data leakage monitoring capabilities the harness provides during execution. Look for held-out set isolation verification ensuring test sets are never used for training or tuning. The harness should detect benchmark material in model training data through multiple techniques: exact match detection via hashing comparing sample signatures against training manifests, fuzzy matching identifying near-duplicates through similarity measures, and n-gram overlap analysis detecting partial matches. Model provenance tracking should query training data sources, check whether evaluation data appears in public datasets, such as Common Crawl or Wikipedia, and analyze model behavior differences on held-out versus previously-seen samples. The harness should alert potential contamination with confidence levels: high for exact matches, medium for high similarity, low for suspicious patterns. Contamination detection logs provide evidence of evaluation integrity for ongoing monitoring as models and benchmarks evolve. \\
\hline

S2F4: Pipeline Orchestration \& Coordination & This capability examines how the harness coordinates evaluation workflow across potentially many interdependent tasks. Look for directed acyclic graph-based dependency management where task dependencies are formalized and parallel execution is enabled for independent tasks. The harness should provide task-type-specific routing with specialized handlers for question-answering, summarization, classification, generation, and custom task types. Protocol variation support should include zero-shot, few-shot with example selection and formatting, chain-of-thought with reasoning prompts, and variation testing for robustness. Task batching should group work efficiently with configurable batch sizes and dynamic batching based on input length. Dynamic load balancing should distribute execution across compute resources with load-aware scheduling. Complete execution traces showing task order, dependencies, and completion status enable debugging and reproducibility. \\
\hline

S2F5: Inference Execution & This capability evaluates the harness's core execution engine for running models on datasets. Look for efficient batching strategies with dynamic sizing adapting to input length and memory or static batching for uniform inputs. Rate limiting should ensure API compliance through requests-per-second limits, concurrent request controls, and exponential backoff on errors. Timeout handling should provide per-sample, per-batch, and global timeouts with graceful degradation. The harness should capture comprehensive generation metadata including latency measurements, token counts for cost tracking, API errors and retry counts, and truncation events. Model constraint respect should handle context window limits through truncation or splitting, verify modality support, and manage API versioning and deprecation. Generation parameters, such as temperature, top-p, and sampling strategy should be configurable and properly applied. \\
\hline

S2F6: Error Recovery \& Resilience & This capability assesses the harness's capabilities for handling execution disruptions to ensure evaluation completion. Look for automatic retry logic with exponential backoff progressively increasing wait times, configurable maximum retry attempts, and timeout specifications at multiple levels. Circuit breaker patterns should detect consecutive failure patterns and pause execution with alerts rather than continuing to waste resources. Graceful degradation should allow evaluation to continue on remaining samples when individual samples fail, marking failures for investigation. The harness should capture detailed error diagnostics including error type classification distinguishing API errors from timeouts from parsing errors, error frequency patterns, failure correlation with input characteristics, and stack traces with execution context. Error classification should distinguish recoverable transient errors from unrecoverable permanent errors, applying appropriate recovery strategies to each. \\
\hline

S2F7: Execution Checkpointing & This capability examines what state persistence capabilities the harness provides for recovery from interruptions. Look for configurable checkpoint frequency options including sample-based intervals, time-based intervals, and on-demand triggering. The harness should preserve complete execution state including results generated so far, random number generator state for reproducibility, data cursor position for resumption, partial batch state, and execution metadata. Incremental resumption should restart from the last checkpoint, skip completed work, and continue with remaining samples without recomputation. The harness should maintain reproducibility ensuring resumed runs produce identical results to uninterrupted runs given the same seed. Checkpoint storage should include compression options and integrity verification through checksums. The capability is particularly valuable for large datasets, expensive models, and unreliable infrastructure, such as spot instances. \\
\hline

\multicolumn{2}{l}{\textbf{STAGE 3: MEASURE -- Scoring, Statistical Validation \& Comparative Analysis (7 capabilities)}} \\
\hline

S3F1: Output Format Validation & This capability assesses the harness's capabilities for validating model outputs against benchmark-specified format requirements before measurement. Look for schema validation checking structured outputs, such as JSON for correct structure, required fields, and type constraints. Format consistency verification should ensure uniform structure across all outputs. Range validation should confirm numerical values fall within expected bounds and categorical values come from valid sets. The harness should provide detailed diagnostics for parsing failures that prevent measurement. Look for capabilities to distinguish format compliance failures where models produced meaningful responses in wrong formats from construct measurement failures where models genuinely do not understand tasks. The validation process itself should be examined for potential bias: the harness should support testing whether strict parsing systematically penalizes capable models with minor format deviations. Documentation should clarify when format requirements capture the target construct versus formatting compliance capability. \\
\hline

S3F2: Measurement Protocol Execution & This capability evaluates how the harness executes the configured measurement protocol with support for diverse modalities. For human judgment, look for crowdsourcing platform integration, rater qualification systems with screening and training, attention checks through gold standards and timing analysis, inter-rater reliability computation, and demographic documentation. For LLM-as-judge, look for judge model selection interfaces, prompt design with structured rubrics, chain-of-thought reasoning prompts, confidence calibration, and bias mitigation through blinding and randomization. For algorithmic metrics, look for diverse scoring functions with configurable parameters and threshold optimization. Ensemble approaches should combine modalities through voting schemes and confidence weighting with disagreement resolution. The harness should enforce a standardized measurement schema where all modalities return scores, uncertainty estimates, rationales, and metadata. Subjective judgment aggregation should preserve variability rather than forcing premature single-point estimates. \\
\hline

S3F3: Output Result Validation & This capability evaluates what sanity check capabilities the harness provides on evaluation results before reporting. Look for baseline ordering validation checking that random baselines are worst, strong baselines are competitive, and human performance is near ceiling. Metric pattern checking should detect zero variance indicating uniform predictions, perfect scores suggesting potential leakage, and inverse correlations between metrics that should agree. Score range validation should verify probabilities fall between zero and one, counts are non-negative, and bounded metrics stay within bounds. Statistical assumption verification should include normality tests for parametric methods, independence checks for correlation tests, and homoscedasticity tests for analysis of variance. The harness should support cross-validation against prior work by comparing performance levels with published benchmarks. Implausible result detection should flag superhuman performance on difficult tasks, dramatic performance drops on similar tasks, or inconsistent relative performance across metrics. The validation should identify implementation bugs, data quality issues, and measurement failures before they propagate to reports. \\
\hline

S3F4: Statistics Aggregation & This capability examines what comprehensive descriptive statistics capabilities the harness provides for summarizing performance beyond simple averages. Look for central tendency measures including mean, median, and mode, dispersion measures including standard deviation, variance, and interquartile range, and percentile computation including quartiles, deciles, and full distributions. The harness should support score distribution visualization through histograms and kernel density estimates. Distributional property characterization should include skewness revealing asymmetry, kurtosis showing tail weight, modality detection identifying multimodal score distributions, and tail behavior analysis examining outlier prevalence and extreme performance. Aggregate uncertainty computation should include population standard error, confidence intervals on mean scores, and credible intervals on distribution parameters. Temporal aggregation should track evaluation evolution by checkpoint, run, or date. Uncertainties should propagate properly through aggregation calculations. \\
\hline

S3F5: Baseline Comparison & This capability examines how the harness ensures fair comparison conditions when applying measurement protocols to baseline systems. Look for mechanisms enforcing identical data splits across all systems, consistent measurement protocols with the same scoring functions and raters or judges, equal hyperparameter tuning budgets providing comparable optimization effort, and comparable computational resources maintaining similar batch sizes and timeouts. The harness should support diverse baseline types: random baselines using uniform random or majority-class predictions, majority baselines predicting the most frequent class, well-tuned classical methods with optimized hyperparameters, state-of-the-art models from published literature, and human performance from expert annotations or crowd consensus. Result presentation should show target system performance alongside baseline context with difference quantification. This comparative structure enables convergent validity evidence confirming the target outperforms simpler alternatives and discriminant validity assessment positioning performance relative to human capability. \\
\hline

S3F6: Uncertainty Quantification & This capability assesses the harness's capabilities for computing uncertainty estimates and performing statistical testing following pre-specified analysis plans. Look for confidence interval construction using bootstrap resampling with configurable iterations and methods, such as percentile or bias-corrected accelerated, or parametric approaches when distributional assumptions are verified through normality tests. Statistical significance testing should include permutation tests for non-parametric comparisons, paired and independent t-tests, and Wilcoxon signed-rank tests as non-parametric alternatives. Pre-specified multiple comparison corrections should be available including Bonferroni adjustment and Benjamini-Hochberg false discovery rate control. Effect size computation should include standardized mean differences, bias-corrected versions for small samples, and confidence intervals on effect sizes themselves. The harness should quantify multiple uncertainty sources separately: sampling variability from finite dataset size, measurement variability from rater disagreement or stochastic sampling, and model variability from random seed differences. Uncertainty decomposition should separate these components, and uncertainty propagation should carry uncertainties through aggregations and transformations using Taylor approximation or Monte Carlo methods. The harness should support multi-seed evaluation to capture run-to-run variability and distinguish statistical significance measured by p-values from practical significance measured by effect size thresholds. All metrics should be reported with explicit uncertainties to enable confidence in comparative claims. \\
\hline

S3F7: Fairness Assessment & This capability evaluates the harness's capabilities for assessing measurement consistency across demographic groups and protected attributes. Look for support to compute performance metrics separately for subgroups defined by gender, race, ethnicity, age, language, geographic region, and socioeconomic indicators. The harness should provide fairness criterion testing including demographic parity testing equal selection rates, equalized odds testing equal true positive and false positive rates, and predictive parity testing equal precision across groups. Intersectional analysis should examine attribute combinations, such as race crossed with gender since marginal analysis alone is insufficient. Statistical significance testing within subgroups should account for multiple comparisons and reduced sample sizes. For human raters, look for rater demographic documentation capabilities, rater agreement pattern analysis for demographic bias, and rater-ratee demographic concordance analysis. Results should clarify whether performance differences reflect construct measurement versus measurement bias. \\
\hline

\multicolumn{2}{l}{\textbf{STAGE 4: INTERPRET \& RELEASE -- Analysis, Reporting \& Reproducibility (6 capabilities)}} \\
\hline

S4F1: Stratified Performance Analysis & This capability assesses the harness's capabilities for understanding performance variations through flexible stratification. Look for support to stratify by task characteristics including task type, domain, difficulty level, input length, and output length requirements. Metadata-based stratification should slice by source, language, creation date, and annotation quality. Input property stratification should analyze by lexical diversity, syntactic complexity, and semantic ambiguity. The harness should support hierarchical slicing with nested stratifications that drill down progressively, for example from domain to task type to difficulty level. Statistical significance testing should identify performance disparities across subgroups while accounting for multiple comparisons. The harness should facilitate linking performance variations back to benchmark construct definitions to distinguish expected differential performance from unexpected patterns suggesting confounds. Multi-objective tradeoff analysis through Pareto frontiers should identify efficient models on dimensions, such as accuracy versus latency or performance versus cost. \\
\hline

S4F2: Leaderboard Management & This capability assesses the harness's capabilities for systematic comparison across multiple models through rigorous frameworks. Look for pairwise statistical testing comparing model performance via paired or independent t-tests, Wilcoxon signed-rank tests, and permutation tests, always with pre-specified multiple comparison corrections. Effect size computation should include standardized mean differences, bias-corrected versions, and confidence intervals on effect sizes quantifying practical significance. Multi-objective comparison should analyze tradeoffs through Pareto frontiers identifying accuracy versus latency efficient solutions, performance versus cost frontiers, and accuracy versus fairness tradeoff curves. Probabilistic ranking systems should provide robust model ordering through approaches, such as Elo ratings, Bayesian skill ratings, or Bradley-Terry pairwise comparison models. Comparative leaderboards should aggregate results with time-stamped rankings tracking evolution, confidence intervals on rankings reflecting uncertainty, and filtering by benchmark, task, or domain. Multi-metric aggregation should support explicit weighting with uncertainty propagation and transparent handling of Pareto-incomparable models. \\
\hline

S4F3: Failure Mode Analysis & This capability evaluates the harness's capabilities for conducting qualitative and quantitative analysis of failure patterns. Look for automated error clustering using semantic similarity to embed and group failures through clustering algorithms. Error type categorization should distinguish factual errors from reasoning errors from format errors from refusal errors and custom categories. Outlier detection should identify unusual failures not fitting common patterns. Root cause analysis should correlate failures with task types and input characteristics, such as length and complexity. The harness should extract representative samples from each error cluster, generate human-readable explanations, and quantify failure frequencies. Look for capabilities to link failures back to construct definitions, distinguishing construct-related failures where models genuinely do not understand the phenomenon from confounds where failures are driven by non-targeted factors. When failures correlate with confounding factors, this should be documented as evidence potentially limiting construct validity. \\
\hline

S4F4: Interactive Exploration & This capability examines what interactive interface capabilities the harness provides for hands-on result investigation. Look for a sample browser with filtering by score ranges, error types, metadata attributes, and model predictions. Full-text search should enable finding samples across inputs, outputs, rationales, and metadata. Drill-down navigation should move from aggregate metrics to individual samples while preserving context. Custom metric computation should allow on-the-fly definition and application to filtered samples with real-time visualization. Comparative sample views should show side-by-side output comparison across models with difference highlighting and agreement-disagreement analysis. Annotation interfaces should support tagging samples with categories, adding notes, and flagging for investigation. Jupyter notebook integration should enable exporting filtered samples, programmatic analysis, hypothesis testing, and publication-quality visualization. Interactive capabilities enable validation of automated findings, discovery of unexpected patterns, and ad-hoc hypothesis testing. \\
\hline

S4F5: Comparative Report Generation & This capability evaluates the harness's capabilities for generating reports tailored to diverse stakeholder needs. Look for executive summary generation emphasizing key findings and critical limitations with actionable recommendations. Technical deep-dive reports should provide complete statistical details including metric tables with uncertainties, significance test results with p-values and effect sizes, fairness analyses with subgroup breakdowns, failure mode analyses with examples, and stratified performance. Compliance documentation should include audit trails, provenance chains, validation evidence, and fairness certifications for regulated domains. Research paper formats should structure results for publication with methods sections, results sections with tables and figures, and validity discussions. Rich visualizations should include confusion matrices, ROC curves with confidence intervals, performance distributions, fairness heatmaps, failure clusters, and tradeoff scatter plots. The harness should integrate benchmark validity evidence prominently in reports including construct definitions, sampling rationales, known limitations, and explicit validity scope statements with caveats about construct validity boundaries. \\
\hline

S4F6: Artifact Management & This capability assesses the harness's capabilities for capturing execution metadata and packaging complete evaluation artifacts for exact reproduction. Look for comprehensive runtime metadata capture including execution timestamps with start, end, and duration, complete configuration snapshots with all parameters serialized and defaults explicit, model identification with IDs, version hashes, checkpoint locations, and API versions, and environment specifications documenting OS version, Python version, complete library dependency trees, and GPU/CPU specifications. Data provenance tracking should record cryptographic hashes of datasets, split specifications, preprocessing pipelines, and normalization statistics. Execution logs should capture errors with stack traces, warnings and informational messages, performance metrics, and intermediate results. Git integration should track configuration changes and link results to code versions via commit hashes. For external reproducibility, look for complete code packaging including harness implementation, measurement protocols, and statistical analysis scripts with pinned dependencies through requirements files, environment specifications, or Docker containers. Configuration files should capture all parameters, command-line arguments, and environment variables. Dataset accessibility should be provided through download scripts, API access details, preprocessing code, and checksums for verification. Model identifiers should specify sources with complete version information. Random seeds should document all stochastic process seeds. Complete results with provenance should include raw outputs, intermediate results, final metrics, and execution logs bundled together. Artifact storage should include compression options and integrity verification through checksums. The capability should enable independent verification by external researchers, support scientific peer review, facilitate regulatory audit, and provide clear documentation of any aspects preventing perfect reproduction, such as proprietary models or restricted datasets. \\
\hline

\end{longtable}
\restoregeometry